\section{Preliminaries}\label{sec:prelim}

\stitle{Graph notation.}
We work with a simple undirected graph $G=(V,E)$.
%
For a vertex $v$, $N_G(v)=\{u\in V\mid (u,v)\in E\}$ is its neighbor set and $\deg_G(v)=|N_G(v)|$ is its degree.
%
For $X\subseteq V$, $G[X]$ is the subgraph induced by $X$.
%
The implementation stores $G$ in CSR form with an \texttt{offsets} array and an \texttt{edges} array.

\begin{definition}[$k$-Clique]\label{def:k-clique}
A \emph{$k$-clique} is a vertex set $C\subseteq V$ with $|C|=k$ such that every two distinct vertices in $C$ are adjacent.
\end{definition}

\begin{definition}[$s$-Clique Support]\label{def:s-supp}
Given $s\ge2$, the \emph{$s$-clique support} of a vertex $v$ in $G$ is
\[
  \sdeg(v) \;=\; \bigl|\{S\subseteq V\mid v\in S,\ |S|=s,\ S\text{ is a clique in }G\}\bigr|.
\]
\end{definition}

$\sdeg(v)$ is evaluated inside the current induced subgraph when a peel level is being defined; the subscript is dropped when the graph is clear.

\begin{definition}[$k$-$(1,s)$-Nucleus]\label{def:nucleus}
Let $X\subseteq V$.
%
The induced subgraph $G[X]$ is a \emph{$k$-$(1,s)$-nucleus} if:
\begin{enumerate}[leftmargin=*]
  \item every vertex in $X$ belongs to at least $k$ $s$-cliques contained in $G[X]$;
  \item $G[X]$ is \emph{$s$-connected}: for any two vertices $u,v\in X$, there is a sequence $u=v_1,\ldots,v_t=v$ such that each pair $\{v_i,v_{i+1}\}$ is contained in an $s$-clique of $G[X]$;
  \item $G[X]$ is maximal with respect to adding vertices while preserving the first two conditions.
\end{enumerate}
\end{definition}

\begin{definition}[Core Number]\label{def:core}
Given $s\ge2$, the \emph{$(1,s)$-nucleus core number} of a vertex $v$, denoted $\kappa_s(v)$, is the largest $k$ for which $v$ belongs to some $k$-$(1,s)$-nucleus of $G$.
\end{definition}

\stitle{The Clique Path Index.}\label{sec:prelim-cpi}
The Clique Path Index, or \cpi{}, compactly represents $s$-cliques without listing them one by one~\cite{NuclearCD}.

\begin{definition}[\cpi{} Path]\label{def:path}
A \emph{\cpi{} path} is a pair $\TPath=(\Vh(\TPath),\Vp(\TPath))$ of disjoint vertex sets such that $\Vh(\TPath)\cup\Vp(\TPath)$ is a clique in $G$ and $|\Vh(\TPath)|\le s\le|\Vh(\TPath)|+|\Vp(\TPath)|$.
%
The set $\Vh(\TPath)$ is the \emph{hold set}: every encoded $s$-clique contains all of these vertices.
%
The set $\Vp(\TPath)$ is the \emph{pivot set}: an encoded $s$-clique chooses vertices from this pool.
%
The \emph{pivot requirement} of $\TPath$ is $\eta_{\TPath}=s-|\Vh(\TPath)|$, the number of pivot vertices needed to complete the hold set into an $s$-clique.
%
The set of $s$-cliques encoded by $\TPath$ is
\[
  \mathcal{C}(\TPath)
  =
  \bigl\{\Vh(\TPath)\cup P'\mid P'\subseteq\Vp(\TPath),\ |P'|=\eta_{\TPath}\bigr\}.
\]
Its cardinality is $\binom{|\Vp(\TPath)|}{\eta_{\TPath}}$.
\end{definition}

The \cpi{} is a set $\mathcal{P}$ of paths such that every $s$-clique is encoded by exactly one path.
%
We write
\[
  \Sig=\sum_{\TPath\in\mathcal{P}}\bigl(|\Vh(\TPath)|+|\Vp(\TPath)|\bigr)
\]
for the vertex--path incidence count.
%
We also write $\deg_{\Sig}(v)=|\{\TPath\mid v\in\Vh(\TPath)\cup\Vp(\TPath)\}|$, $d_{\max}=\max_v\deg_{\Sig}(v)$, and $\kmax=\max_v\sdeg(v)$.
%
The \cpi{} is built by a degeneracy-ordered Bron--Kerbosch search with Tomita pivoting~\cite{NuclearCD,BronKerbosch,eppstein,tomita}.

\begin{lemma}[Support Formula~\cite{NuclearCD}]\label{lem:supp}
For every $v\in V$,
\[
  \sdeg(v)
  =
  \sum_{\TPath\,:\,v\in\Vh(\TPath)}\binom{|\Vp(\TPath)|}{\eta_{\TPath}}
  +
  \sum_{\TPath\,:\,v\in\Vp(\TPath)}\binom{|\Vp(\TPath)|-1}{\eta_{\TPath}-1}.
\]
\end{lemma}

\sstitle{Problem statement.}
Given $G$ and $s\ge2$, compute $\kappa_s(v)$ for every $v\in V$; the nucleus hierarchy is constructed afterwards (Section~\ref{sec:buildhier}).
